\chapter{Background and related work}
\label{chap:background}

\section{Browser agents and indirect prompt injection}

An LLM agent places a generative model in a feedback loop of goals,
observations, decisions, tool calls, and results. ReAct made this
reasoning--action interleaving explicit, while Toolformer examined how models
can learn to invoke tools~\cite{yao2023react,schick2023toolformer}. Browser
agents bring the same loop into environments composed of URLs, structured page
representations, interactive elements, and code-capable actions. WebArena and
WorkArena evaluate such agents in executable environments with task-specific
evaluators, including state and answer checks
~\cite{zhou2024webarena,drouin2024workarena}. Together, these benchmarks
highlight two properties relevant here: success is a validated task outcome,
not merely a plausible response, and the agent's observations contain data
outside the operator's control.

This exposure creates a security problem because instructions and data share
the same model context. Direct prompt-injection attacks alter the user's
prompt; indirect attacks instead arrive through retrieved content that the
application expected to treat as data.
Greshake et al. demonstrated indirect prompt-injection attacks in synthetic
LLM-integrated applications and through local web content, including data
exfiltration and worm-like propagation~\cite{greshake2023indirect}. In a
browser setting, an attacker can place equivalent instructions in a page,
user-generated record, link label, or stored result, where the agent encounters
them while acting with the user's delegated authority. For example, injected
issue content can divert an agent asked to comment on an issue towards account
settings, with the attacker goal of making the benchmark user's profile public
~\cite{evtimov2025wasp}. The page supplies the hostile instruction, but the
resulting navigation occurs within the benchmark user's authenticated sandbox
session.

Indirect prompt injection therefore differs from conventional parser
injection. A browser agent is deliberately expected to interpret previously
unseen prose, so a fixed grammar cannot by itself establish which source has
authority. Delimiters can mark the distinction between trusted instructions
and untrusted input, but page text can imitate them, leaving the model to decide
which instruction deserves authority. Formal analyses and benchmark studies
reinforce this point: models do not reliably distinguish the instructions they
should follow from external content they should treat as data
~\cite{liu2024formalizing,yi2025bipia}.

Nor is the threat confined to the current page. Retrieved content can influence
subsequent model responses and API calls, while persistent memory can
reintroduce hostile content in later sessions~\cite{greshake2023indirect}.
Multimodal agents add a further channel: image-hijack work shows that specially
optimised adversarial images can carry the behavioural effect of a chosen text
prompt and redirect a vision-language model's behaviour
~\cite{bailey2024imagehijacks}. A filter over the Document Object Model (DOM),
the browser's structured representation of a page, cannot remove content
rendered only in pixels. Visual observations therefore form a distinct threat
surface that requires its own controls; text sanitisation does not automatically
cover it.

\subsection{What an agent-security evaluation must measure}

These threat surfaces also make the unit of evaluation important. A detector
labels a piece of text; an action monitor labels a proposed or executed effect;
and an end-to-end oracle checks whether the user or attacker changed the
environment as intended. High detector recall can still coexist with a
compromised trajectory if the model is merely warned or a dangerous effect
bypasses the monitored channel. Conversely, low terminal attack success can
describe an agent that was diverted but lacked the capability or remaining
steps to complete the attack.

AgentDojo operationalises this distinction by pairing realistic user tasks with
attacker goals in executable banking, messaging, travel, and workspace
environments~\cite{debenedetti2024agentdojo}. Its native oracles score both
goals through task-specific checks that may inspect environment state and model
output. WASP makes a related distinction for browser agents, using an LLM judge
to assess whether a trajectory has been diverted before the terminal attacker
goal is achieved~\cite{evtimov2025wasp}. The benchmarks are therefore
complementary: AgentDojo provides controlled tool-agent environments, whereas
WASP adds browser-native observations, navigation, and re-planning. Their
absolute rates cannot be pooled, however, because their tasks, action spaces,
and oracles differ. ST-WebAgentBench broadens the picture further by studying
policy compliance through Completion Under Policy and Risk Ratio
~\cite{levy2026stwebagentbench}. These measures must be read together because
detection, risky execution, policy compliance, and final-state success reveal
different parts of system behaviour.

Accordingly, this dissertation treats security and utility as a joint outcome.
It reports native user completion, native attacker success, useful-and-secure
completion, attacker-directed attempts and executions, and recorded
enforcement events alongside trajectory outcomes, while separating scripted
controls from actual human intervention. This combination guards against two
misleading conclusions: calling an inactive agent secure, or calling a capable
but compromised agent useful.

\section{Defence design space}

\subsection{Model-facing defences}

One family of defences tries to help the planner recognise which content
deserves instructional authority. BIPIA evaluates reminders, boundary
awareness, in-context examples, and training across indirect-injection tasks
~\cite{yi2025bipia}. Spotlighting takes a complementary approach, using
delimiting, datamarking, or encoding to keep the distinction between trusted
instructions and untrusted input salient to the model
~\cite{hines2024spotlighting}.

Other approaches try to make this separation part of model behaviour. StruQ
uses structured instruction tuning to distinguish prompt and data channels
~\cite{chen2025struq}, while SecAlign uses preference optimisation to favour
responses that follow the legitimate instruction over the injection
~\cite{chen2025secalign}. Both report greater resistance to the evaluated
attacks with limited utility loss in their respective settings. More
specifically, on its principal AlpacaFarm evaluation, SecAlign reports 0\%
attack success for the evaluated optimisation-free attacks and rates mainly
below 10\% for stronger optimisation-based attacks, with utility similar to its
non-preference-optimised counterparts. A separate out-of-distribution SEP
Ignore evaluation reports 11.0\%, 3.6\%, and 6.5\% attack success when the
injection appears at the start, middle, and end of the input. These findings
remain bound to the paper's models, tasks, and attacks. A later diagnostic shows
why attack rejection and benign utility must be read together. Across three
reasoning benchmarks, benign suffix-task rejection in the evaluated fine-tuned
defences rose from below 10\% to as high as 90\%. The authors link these
failures to position, trigger-token, and topic correlations, suggesting that
trained defences can rely on narrow surface cues without reliably identifying
malicious intent~\cite{li2026surfaceheuristics}.

Model-facing checks can also operate during inference. Task Shield asks whether
each instruction and proposed tool call contributes to the user's stated goal
~\cite{jia2025taskshield}. Under the one-trial-per-task Important
Instructions attack condition on AgentDojo with GPT-4o, its authors report
2.07\% attack success and 69.79\% utility under attack. MELON instead reruns
the agent trajectory with the user prompt masked and treats similarity between
the original and counterfactual tool calls as evidence of hijacking. The paper
notes that response-based attacks without tool calls can evade this comparison
~\cite{zhu2025melon}. Although their mechanisms differ, both examine whether
planned behaviour remains dependent on the authorised task and report
AgentDojo trade-offs within their own configurations. Sentinel additionally
tests transfer through WASP's browser-native loop. Even within AgentDojo,
however, differences in models, attack construction, scaffolds, defence
mechanisms, and metric definitions make headline rates contextual; they do not
provide a controlled ranking against Sentinel.

Detection offers another inspectable intervention before processed content
reaches the downstream LLM. Detect-and-Remove separates document-level attack
identification from sentence-level removal or generative extraction followed
by longest-common-substring removal~\cite{chen2025detectremove}. The approach
therefore depends on detecting the attack and identifying what can be removed
without losing information needed to complete the task. Distribution shift
adds a broader challenge: no fixed collection of examples can cover every
wording, encoding, page context, or adaptive strategy. Indeed, white-box
optimisation against specific defences has outperformed non-adaptive attacks in
the evaluated settings~\cite{zhan2025adaptive}. These limitations motivate
Sentinel's layered use of detection: it governs observations while policy is
applied independently before execution. Model-facing guidance is thus backed
by external checks on what the agent observes and can do, keeping execution
authority outside the model.

\subsection{System-level authority}

Where model-facing defences influence planning, system-level work moves selected
security and execution decisions outside unconstrained model reasoning. Some
designs nevertheless retain LLM components within the wider control
architecture. CaMeL, for example, extracts control and data flow from a trusted
query, confines untrusted data to constrained processing, and enforces
capabilities at tool calls~\cite{debenedetti2025camel}. Fides instead applies
information-flow control: confidentiality and integrity labels propagate
through planning, deterministic policies govern tool effects, and information
can be selectively hidden from the planner~\cite{costa2025fides}. Progent
represents privilege as symbolic rules over tool names and arguments, checks
every tool call deterministically, applies privilege narrowing automatically,
and can require approval before an expansion. Its main experiments use
Auto-Approve, so they do not measure human-reviewed expansion
~\cite{shi2025progent}. Building on Fides, PRUDENTIA makes planning policy-aware
and evaluates autonomy through task completion at intervention budgets and
trace-derived oversight load under assumed approval
~\cite{kolluri2026prudentia}.

Other systems place trust in different abstractions. AgentSpec provides a
lightweight language of triggers, predicates, and runtime enforcement actions
~\cite{wang2026agentspec}. DRIFT derives a minimal function trajectory and
parameter checklist from the trusted query, validates deviations, and masks
conflicting instructions before they persist in memory~\cite{li2025drift}.
IPIGuard plans an initial tool-dependency graph before exposure to external data
and constrains later execution to its traversal, while allowing read-only query
nodes to be added during execution~\cite{an2025ipiguard}. VIGIL takes a
different route: it permits speculative reasoning over the tool stream, then
verifies a proposed effect against the trusted intent before commitment
~\cite{lin2026vigil}. AttriGuard instead tests whether a proposed tool call
survives teacher-forced shadow replay under control-attenuated external
observations~\cite{he2026attriguard}. Their reported attack-success rates are
not directly comparable, since each system changes a different part of the
agent stack and is evaluated under different policies, models, and tasks.
Moreover, each assurance remains conditional on the system's trusted
abstraction and expressible policy. \Cref{tab:system-comparison} therefore asks
the more durable architectural question: what does each system trust, and where
is that trust enforced?

\begin{table}[!t]
\centering
\caption{Representative system-level approaches.}
\label{tab:system-comparison}
\footnotesize
\renewcommand{\arraystretch}{1.08}
\setlength{\tabcolsep}{3.5pt}
\begin{tabularx}{\textwidth}{@{}L{0.16\textwidth}L{0.22\textwidth}L{0.22\textwidth}Y@{}}
\toprule
System & Trusted abstraction & Enforcement boundary & Evidence and assurance \\
\midrule
CaMeL~\cite{debenedetti2025camel} &
Trusted-query program and capabilities &
Structured flow and tool calls &
By-design guarantees for supported flows; AgentDojo. \\
\addlinespace[2.25pt]

Fides\newline\cite{costa2025fides} &
Security labels and flow policy &
Planner exposure and tools &
Formal IFC; AgentDojo. \\
\addlinespace[2.25pt]

Progent\newline\cite{shi2025progent} &
Generated symbolic tool and argument privileges &
Every call and privilege expansion &
Deterministic confinement; AgentDojo and ASB. \\
\addlinespace[2.25pt]

PRUDENTIA\newline\cite{kolluri2026prudentia} &
Fides labels and policy-aware planning &
Planning, exposure, endorsement, and tools &
Inherited deterministic IFC; AgentDojo and WASP under assumed approval. \\
\addlinespace[2.25pt]

AgentSpec~\cite{wang2026agentspec} &
Trigger--predicate--enforcement rules &
Configured runtime events &
Empirical evaluation in code, embodied, and driving domains. \\
\addlinespace[2.25pt]

AttriGuard~\cite{he2026attriguard} &
Survival under control-attenuated observations &
Shadow replay before tool execution &
Static: four LLMs and two benchmarks; adaptive: two backbones on a 61-case
AgentDojo subset. \\
\addlinespace[2.25pt]

\systemname &
Trusted request, policy, live context, and narrow server-owned authority &
Wrapped model-input regions and pre-execution actions &
Contract tests; AgentDojo and matched WASP outcomes. \\
\bottomrule
\end{tabularx}
\end{table}

The table shows differences in coverage, not a universal ordering. Progent and
AttriGuard mediate the action boundary; Sentinel also
governs wrapped planner observations. Sentinel nevertheless preserves
open-ended replanning, so policy can reject prohibited effects but cannot prove
that an allowed click advances the operator's goal. Its claim is therefore
limited to the evaluated observation, action, and authority boundaries, not
universal superiority.

\section{Protection principles for agent runtimes}

These system-level approaches connect to established protection principles.
The reference-monitor concept requires a mechanism to be invoked on every
protected operation, isolated from the subject it constrains, and small enough
to analyse~\cite{anderson1972computersecurity}. Complete mediation, fail-safe
defaults, and least privilege provide corresponding design tests
~\cite{saltzer1975protection}. A language model is a poor monitor for its own
effects because hostile input can influence its internal state. A
wrapper at the executor can instead satisfy the structural requirement that
every action in its governed tool path crosses one decision point, even though
natural-language risk classification remains empirical.

Information-flow control provides a complementary perspective by asking where
data originated and which sinks it may reach. Denning's lattice model
formalised label-constrained flow~\cite{denning1976lattice}. Browser-agent
values are less structured, however: page content becomes model context, and a
click can disclose information without an explicit network API. At each action
assessment, Sentinel therefore derives a bounded dynamic model. Wrapped browser
observations carry untrusted provenance, detections add injection taint,
recognised values add sensitive categories, and actions expose capabilities
such as navigation, page write, file access, and code execution. These labels
give the policy engine concrete source--sink facts to enforce at execution
time. They remain operational controls, however, and do not prove that every
information flow is isolated.

Finally, runtime enforcement needs a precise authority protocol. Effective
human-in-the-loop security depends on communication, processing, application,
and resulting behaviour~\cite{cranor2008humanloop}; automation research
likewise distinguishes information acquisition, information analysis, decision
and action selection, and action implementation
~\cite{parasuraman2000automation}. Drawing on this separation, Sentinel maps
hard refusal, narrow approval, takeover, and observation review to distinct
server outcomes. This structure prevents a generic confirmation from weakening
a hard rule. Whether the interface helps people make sound choices remains a
separate human-subject question, particularly because warning design affects
attention and response~\cite{sunshine2009cryingwolf,akhawe2013warningland}.

\section{Research gap}

Taken together, prior work spans learned separation, detection, structured
control, information flow, and runtime rules. The question motivating this
dissertation is therefore not whether one defence can eliminate prompt
injection. It is whether an independent layer placed between a browser agent
and its privileged executor can govern both the observations that shape
planning and the actions that produce effects, while retaining useful
replanning and reducing attacker-serving execution. Sentinel addresses this
narrower gap by combining model-input
mediation with pre-execution policy and server-owned authority. Its evidence is
deliberately divided: AgentDojo examines the enforcement mechanism under
controlled tool-agent conditions; WASP tests transfer to a browser-native
observation and navigation loop; and contract tests with deterministic
walkthroughs provide implementation evidence for the operator-authority
workflow. Together, these strands ask whether evaluated policy violations can
be reduced without eliminating useful autonomy, rather than claiming a general
solution to prompt injection or evidence about human behaviour.
